%
% Front matter
%

\begin{frontmatter}

\title{
  T2K \TotalPot-POT 3-Flavour Electron-Antineutrino Appearance Analysis
}

\vspace*{0.7cm}

\author[RAL,LIV]  {Costas Andreopoulos}
%\author[ETHZ]     {Fatih Bay}
\author[RAL,OXF]  {Thomas Dealtry}
\author[RAL,WAR]  {Steve Dennis}
\author[IFIC]     {Lorena Escudero S\'anchez}
\author[LANCS,WAR]{\\Nick Grant}
\author[ETHZ]     {Davide Sgalaberna}
\author[RAL,OXF]  {Raj Shah}
\author[RAL,OXF]  {and Alfons Weber}

\address[RAL]   {Science \& Technology Facilities Council (STFC), Rutherford Appleton Laboratory,\\ Harwell Oxford Campus, OX11 0QX, UK}
\address[LIV]   {University of Liverpool, Physics Department,\\ Liverpool, L69 7ZE, UK} 
\address[ETHZ]  {ETH Zurich, Institute for Particle Physics,\\ 8092 Zurich, Switzerland}
\address[OXF]   {University of Oxford, Subdepartment of Particle Physics,\\ Oxford, OX1 3RH, UK}
\address[WAR]   {University of Warwick, Physics Department,\\ Coventry, CV4 7AL, UK}
\address[IFIC]  {CSIC and University of Valencia, Instituto de Fisica Corpuscular, \\ E-46071 Valencia, Spain}
\address[LANCS] {Lancaster University, Physics Department,\\ Lancaster, LA1 4YB, UK}

%
% Abstract
%

\begin{abstract}

We report the results of the
3-flavour $\bar{\nu}_{e}$ appearance analysis, performed by the {\em VALOR} group,
on the combined Run 5+6 dataset, using only data taken in antineutrino beam mode.
This analysis adapts the work presented in Refs. 
\cite{OurRun1DisappAnalysis},
\cite{OurRun12DisappAnalysis}, 
\cite{OurRun12DisappAnalysis3Flavour},
\cite{OurRun123DisappAnalysis3Flavour},
\cite{OurRun1234aDisappAnalysis3Flavour},
\cite{OurRun1234DisappAnalysis3Flavour},
\cite{OurRun123JointOscillationAnalysis},
\cite{OurRun1234JointOscillationAnalysis} and
\cite{OurRun56AntiNumuDisappearanceAnalysis3Flavour}
to this dataset.
The Run 5+6 dataset corresponds to an integrated J-PARC neutrino beam
exposure of \TotalPot{} POT.

This analysis predicts 
\NskOscExpectation{} $\pm$ \NskOscExpectationError{} (syst)
single $e$-like ring events in Super K  considering $\bar{\nu}_e$ appearance
and  
\NskObserved{} 
were observed. 

%The observed deficit has a strong energy dependence: the ratio of the number of observed events to the 
%number expected under the no-oscillation hypothesis is
%$\sim$\NskRatioLowEnergy\% $<$ 0.5 GeV, 
%$\sim$\NskRatioMidEnergy\% between 0.5 and 1 GeV and 
%$\sim$\NskRatioHighEnergy\% $>$ 1 GeV. 

%In order to be able to investigate whether the oscillation parameters for $\bar{\nu}_{\mu}$ disappearance are different
%from those for $\nu_{\mu}$ disappearance, two separate sets of oscillation parameters are used.
%One set is applied to the neutrino interaction modes,
%while the other set is applied to the antineutrino modes.

A $\bar{\nu}_{e}$-appearance analysis was performed 
in a framework of 3-flavour oscillations including matter effects in constant-density matter.

The observed reconstructed energy spectrum of single $e$-like ring events was compared
to the expected energy spectrum for different signal hypotheses: absence and presence of $\bar{\nu}_e$
appearance consistent with the current best knowledge of the PMNS matrix,
i.e. T2K+reactor joint fit with the Run1-4 data set.
%The observed reconstructed energy spectrum of single $e$-like ring events was fitted, 
%and separate fits were made for the normal and the inverted mass hierarchies. 

A new not physical parameter, $\beta$, was introduced to define the absence or the presence 
of $\bar{\nu}_e$ appearance.
All the \NSyst{} systematic parameters, which include the oscillation parameters as well as the mass hierarchy, were marginalized,
fixing $\beta$ either to 0 or 1.
The p-value given the null $\bar{\nu}_e$ appearance hypothesis was quoted
performing many toy experiments.
The quoted rate-only p-value is \PvalueRate, while the rate+shape p-value is \PvalueShape.

%In these fits, 
%$\sin^{2}\overline{\theta}_{23}$ and 
%either 
%$|\overline{\Delta m}^{2}_{32}|$ 
%(normal hierarchy) or $|\overline{\Delta m}^{2}_{31}|$ (inverted hierarchy) 
%were allowed to float, while all other oscillation  parameters in both sets were fixed to T2K run1-4 best-fit values.
%All \NSyst{} systematic parameters considered in this analysis were also allowed to float in the fits.

%The 3-flavour $\nu_{\mu}$-disappearance fit of the combined Run 5+6 dataset for the normal mass hierarchy
%gives $\sin^{2}\overline{\theta}_{23}$ = \BestFitMixingNormal{} and $|\overline{\Delta m}^{2}_{32}|$ = \BestFitDmsqNormal{} 
%eV$^{2}$/c$^{4}$ as the best-fit values (p-value = \PValueNormal). Assuming the inverted mass hierarchy, the fit yields 
%$\sin^{2}\overline{\theta}_{23}$ = \BestFitMixingInverted{} and $|\overline{\Delta m}^{2}_{31}|$ = \BestFitDmsqInverted{} 
%eV$^{2}$/c$^{4}$ as the best-fit values (p-value = \PValueInverted). The p-value for the normal hierarchy best-fit values 
%given a null hypothesis of the normal hierarchy best-fit values of the equivalent parameters from the official Run 1-4 neutrino 
%mode analysis is 0.74; thus the results of this analysis considered together with the Run 1-4 neutrino mode results are 
%consistent with CPT being conserved. Confidence regions in the parameter space $\sin^{2}\overline{\theta}_{23}$ - $|\overline{\Delta m}^{2}_{32}|$ 
%(normal hierarchy) and $\sin^{2}\overline{\theta}_{23}$ - $|\overline{\Delta m}^{2}_{31}|$ (inverted hierarchy) 
%%~ are shown using both the constant-$\Delta\chi^2$ method and the method of Feldman and Cousins.
%are shown using the constant-$\Delta\chi^2$ method.

\end{abstract}

\end{frontmatter}
